\chapter{Ուսումնական ուղի}
\label{ch:learning}

\section{Նախնական պահանջներ}
\textbf{Ուսանող 1}
\begin{itemize}
  \item Թվային սխեմաներ (կոմբինացիոն/հաջորդական տրամաբանություն, FSM) — \emph{ինքնաստուգում.} կարող է նախագծել counter և FSM, գծել ժամանակային դիագրամ։
  \item Verilog/SystemVerilog կամ VHDL հիմունքներ — \emph{ինքնաստուգում.} գրել և սիմուլացնել պարզ մոդուլ։ Եթե թույլ է՝ \cite{harris2016}, գլ.~4։
  \item C ծրագրավորում — PS-ի ծրագրի համար։
\end{itemize}
\textbf{Ուսանող 2}
\begin{itemize}
  \item Ծրագրավորում (Python կամ JavaScript/TypeScript) — \emph{ինքնաստուգում.} գրել փոքր REST սերվեր։
  \item Տվյալների բազաներ և SQL — \emph{ինքնաստուգում.} նախագծել 3 աղյուսակով սխեմա և գրել JOIN հարցում։
  \item Հավանականություն և վիճակագրություն (mean, variance, բաշխումներ)։
\end{itemize}

\section{Ինչ սովորել}
\levellegend

\begin{learningpath}[Գիտելիքների քարտեզ — \studentlabel{1}.]
  \learngroup{Թվային նախագծում և վերիֆիկացիա}
  \learn{Սինքրոն նախագծում, FSM, pipeline}{\apply}{գեներատորներ, processing IP}{\cite{harris2016}, գլ.~3--4}
  \learn{Testbench և սիմուլյացիա}{\apply}{բոլոր մոդուլների վերիֆիկացիա}{\cite{ug900}; \cite{chu2018}}
  \learn{LFSR, ֆիքսված ստորակետով թվաբանություն}{\apply}{աղմուկ, scaling, ֆիլտրեր}{\cite{chu2018}}
  \learn{Moving average և FIR ֆիլտրեր}{\understand}{processing IP}{\cite{smith1997}, գլ.~14--15}
  \learngroup{Zynq և AXI}
  \learn{Zynq UltraScale+ ճարտարապետություն (PS/PL)}{\understand}{համակարգի ինտեգրում}{\cite{ug1085}; \cite{zynqbook}}
  \learn{AXI4-Lite slave}{\apply}{կոնֆիգուրացիոն ռեգիստրներ}{\cite{axispec}; \cite{zynqbook}}
  \learn{AXI4-Stream և AXI FIFO}{\apply}{տվյալների ուղի PL$\rightarrow$PS}{\cite{pg080}}
  \learn{AXI DMA, DDR}{\aware}{ընդլայնում}{\cite{ug1085}}
  \learngroup{Իրականացում}
  \learn{Vivado block design, IP packaging}{\apply}{PL-ի ինտեգրում}{\cite{zynqbook}}
  \learn{Timing constraints և timing closure}{\understand}{implementation, հաշվետվություն}{\cite{ug906}}
  \learn{Vitis bare-metal ծրագիր, Ethernet (lwIP)}{\understand}{PS-ի ծրագիր}{\missing{AMD-ի lwIP օրինակ — ստուգել}}
\end{learningpath}

\begin{learningpath}[Գիտելիքների քարտեզ — \studentlabel{2}.]
  \learngroup{Backend և տվյալներ}
  \learn{REST API նախագծում}{\apply}{տվյալների ընդունում, հրամաններ}{\cite{fastapi}}
  \learn{Ռելյացիոն մոդելավորում, ժամանակային շարքեր}{\apply}{Experiment/Sample մոդել}{\cite{postgres}; \cite{kleppmann2017}, գլ.~2--3}
  \learn{Հոսքային տվյալներ (WebSocket կամ polling)}{\understand}{live վիզուալիզացիա}{\cite{kleppmann2017}, գլ.~11}
  \learngroup{Վերլուծություն}
  \learn{Նկարագրական վիճակագրություն}{\apply}{փորձերի համեմատում}{\cite{mckinney2022}}
  \learn{Անոմալիաների հայտնաբերում}{\apply}{մոդելավորման բաղադրիչ}{\cite{chandola2009}}
  \learn{Ազդանշանի ֆիլտրում}{\understand}{FPGA և սերվերի մշակման համեմատություն}{\cite{smith1997}, գլ.~15}
  \learngroup{Համակարգ}
  \learn{Zynq PS/PL-ի բաժանում, AXI ռեգիստրներ}{\aware}{կոնֆիգուրացիոն հրամանների իմաստ}{\cite{zynqbook}}
\end{learningpath}

\section{Գործիքներ և հմտություններ}
\begin{itemize}
  \item \textbf{Vivado / Vitis} (ուսանող 1) — սինթեզ, implementation, PS-ի ծրագիր։
  \item \textbf{Python (FastAPI), PostgreSQL, Web frontend} (ուսանող 2) — առաջարկվող տարբերակ. վերջնական ընտրությունը համաձայնեցվում է ղեկավարի հետ։
  \item \textbf{git} (երկուսն էլ) — ընդհանուր ռեպոզիտորիա, branch-եր, code review։
\end{itemize}
